\begin{corollary}[Interior $C^1$ regularity]
Fix an integer $d\ge3$, let $\Omega\subseteq\R^d$ be a nonempty connected set
and let $u:\R^d\to\R$ be infinity-harmonic on $\Omega$ in the sense of
\noderef{IsInfHarmonicOn}, which in particular makes $\Omega$ open and $u$
continuous on $\Omega$. Then $u$ is continuously differentiable on $\Omega$:
there is a map $g:\Omega\to\R^d$, continuous on $\Omega$, such that $u$ is
differentiable at every $x\in\Omega$ with gradient $g(x)$. Equivalently, since
$\Omega$ is open, $u\in C^1_{\mathrm{loc}}(\Omega)=C^1(\Omega)$.
\end{corollary}
\begin{proof}
By \noderef{IsInfHarmonicOn} the set $\Omega$ is open and $u$ is continuous on
$\Omega$. The gradient of $u$ at a point, when it exists, is unique, being the
Riesz representative of the Fr\'echet derivative there; so it suffices to
produce, for every $x_0\in\Omega$, an open ball $V\ni x_0$ with $V\subseteq
\Omega$ such that $u$ is differentiable at every point of $V$ with a gradient
depending continuously on that point. Indeed, the resulting gradient map is then
well defined on all of $\Omega$ — any two of the local maps agree on the
overlap, both being the gradient of $u$ — it is defined at every point of
$\Omega$, and continuity is a local property, so it is continuous on $\Omega$.
Note that the hypotheses that $\Omega$ is nonempty and connected are not used
below; they are carried because the statement being formalized carries them.

\emph{Step 1: a good ball and the rescaled function.}
Fix $x_0\in\Omega$. Since $\Omega$ is open there is $R>0$ with
$B_R(x_0)\subseteq\Omega$. Put $r=R/4>0$, so
\[
  \overline B_{2r}(x_0)\subseteq B_R(x_0)\subseteq\Omega ,
\]
because $|x-x_0|\le2r=R/2<R$. Let $A=2r>0$, let $T(y)=x_0+Ay$ and put
\[
  v=u\circ T,\qquad v(y)=u(x_0+2ry).
\]
By \noderef{InfHarmonicRescale}, applied with this $\Omega$, this $u$, the point
$x_0$ and the factor $A>0$, the function $v$ is infinity-harmonic on
$T^{-1}(\Omega)$.

\emph{Step 2: $v$ is infinity-harmonic on the unit ball.}
First, $B_1\subseteq T^{-1}(\Omega)$: if $|y|<1$ then
$|T(y)-x_0|=2r|y|<2r$, so $T(y)\in\overline B_{2r}(x_0)\subseteq\Omega$.

Next, infinity-harmonicity passes to open subsets. Let $W\subseteq\Omega'$ with
$W$ open, $\Omega'$ open and $f$ an infinity subsolution on $\Omega'$ in the
sense of \noderef{IsInfSubsolutionOn}. Let $z\in W$ and let $\phi$ be $C^2$ near
$z$ with $f-\phi$ having a local maximum on $W$ at $z$, i.e.\ there is a
neighbourhood $N$ of $z$ with $(f-\phi)(x)\le(f-\phi)(z)$ for all $x\in N\cap W$.
Then $N\cap W$ is itself a neighbourhood of $z$, $W$ being open and containing
$z$, and for every $x\in(N\cap W)\cap\Omega'$ we have $x\in N\cap W$ and hence
$(f-\phi)(x)\le(f-\phi)(z)$; so $f-\phi$ has a local maximum on $\Omega'$ at
$z$, and \noderef{IsInfSubsolutionOn} gives $\Dinf\phi(z)\ge0$. As $W$ is open
and $f$ is continuous on $\Omega'\supseteq W$, hence on $W$, this shows that $f$
is an infinity subsolution on $W$. The same argument with inequalities reversed,
using \noderef{IsInfSupersolutionOn}, shows that an infinity supersolution on
$\Omega'$ is one on $W$. By \noderef{IsInfHarmonicOn} — which is the conjunction
of the two — infinity-harmonicity on $\Omega'$ implies infinity-harmonicity on
every open $W\subseteq\Omega'$.

Applying this with $\Omega'=T^{-1}(\Omega)$, which is open because $T$ is a
homeomorphism of $\R^d$ and $\Omega$ is open, and $W=B_1$, we conclude that $v$
is infinity-harmonic on $B_1$ in the sense of \noderef{IsInfHarmonicOn}.

\emph{Step 3: $v$ is bounded on the unit ball.}
By \noderef{IsInfHarmonicOn} the function $u$ is continuous on $\Omega$, so it is
continuous on the compact set $\overline B_{2r}(x_0)\subseteq\Omega$, which is
closed and bounded in $\R^d$ and hence compact. By the Weierstrass item of the
preamble $u$ attains a maximum and a minimum on it; call them $m^+$ and $m^-$.
Every point of $v(B_1)$ has the form $u(T(y))$ with $T(y)\in\overline
B_{2r}(x_0)$, as in Step 2, so $m^-\le u(T(y))\le m^+$. Thus $v(B_1)$ is bounded
above and bounded below.

\emph{Step 4: the estimate for $v$ and the H\"older continuity of its gradient.}
Let $\alpha$ and $C$ be the constants \noderef{MainEstimate} provides for this
$d$, so $0<\alpha\le1/3$ and $0<C<\infty$. By Steps 2 and 3 the function $v$
satisfies its hypotheses, so there is a map $g_v:\R^d\to\R^d$ such that $v$ is
differentiable at every $y\in B_{1/2}$ with gradient $g_v(y)$, the set
$Q_\alpha(g_v;B_{1/2})$ of H\"older quotients of \noderef{HolderQuotSet} is
bounded above, and the two remaining conclusions hold as well; we use only
these two.

Put $K=[g_v]_{C^{0,\alpha}(B_{1/2})}$, the H\"older seminorm of
\noderef{HolderSeminormOn}, which is the supremum of
$Q_\alpha(g_v;B_{1/2})$. That set is nonempty, because $d\ge3\ge1$ makes
$B_{1/2}\subseteq\R^d$ contain two distinct points, and bounded above, so $K$ is
a real number and dominates every member of the set; also $K\ge0$, every member
being a quotient of nonnegative numbers. By \noderef{HolderQuotSet}, for all
$y,y'\in B_{1/2}$ with $y\ne y'$,
\[
  \frac{|g_v(y)-g_v(y')|}{|y-y'|^\alpha}\le K,
  \qquad\text{hence}\qquad
  |g_v(y)-g_v(y')|\le K|y-y'|^\alpha ,
\]
and the last inequality also holds trivially when $y=y'$. Since $\alpha>0$, the
right-hand side tends to $0$ as $y'\to y$, so $g_v$ is continuous at every
$y\in B_{1/2}$ relative to $B_{1/2}$; as $B_{1/2}$ is open, $g_v$ is continuous
on $B_{1/2}$.

\emph{Step 5: unscaling.}
Put $V=B_r(x_0)$, an open ball with $x_0\in V$ and
$V\subseteq\overline B_{2r}(x_0)\subseteq\Omega$, and define
\[
  S(x)=\frac{x-x_0}{2r},\qquad g(x)=\frac{1}{2r}\,g_v\bigl(S(x)\bigr)
  \qquad(x\in V).
\]
For $x\in V$ we have $|S(x)|=|x-x_0|/(2r)<r/(2r)=1/2$, so $S(x)\in B_{1/2}$ and
$g$ is defined on $V$. Note $T(S(x))=x_0+2r\cdot(x-x_0)/(2r)=x$, so
$v(S(x))=u(x)$, and more generally $v\bigl(S(x)+k\bigr)=u(x+2rk)$ for every
$k\in\R^d$, since $T(S(x)+k)=x+2rk$.

Fix $x\in V$ and put $y=S(x)\in B_{1/2}$. For $h\in\R^d$ write $k=h/(2r)$, so
$h=2rk$ and $|k|=|h|/(2r)$. Then
\[
  u(x+h)-u(x)-\Bigl\langle\frac{1}{2r}g_v(y),h\Bigr\rangle
   =v(y+k)-v(y)-\bigl\langle g_v(y),k\bigr\rangle ,
\]
because $u(x+h)=v(y+k)$, $u(x)=v(y)$ and
$\langle(2r)^{-1}g_v(y),2rk\rangle=\langle g_v(y),k\rangle$. By Step 4 the right
side is $o(|k|)$ as $k\to0$, and $|k|=|h|/(2r)$ with $2r$ a fixed positive
constant, so the left side is $o(|h|)$ as $h\to0$. Hence $u$ is differentiable at
$x$ with gradient $g(x)=(2r)^{-1}g_v(S(x))$.

Finally $g$ is continuous on $V$: it is the composition of the continuous affine
map $S$, which carries $V$ into $B_{1/2}$, with the map $g_v$, continuous on
$B_{1/2}$ by Step 4, followed by multiplication by the constant $(2r)^{-1}$.

So $V=B_r(x_0)$ is an open ball containing $x_0$ and contained in $\Omega$ on
which $u$ is differentiable with continuous gradient. Since $x_0\in\Omega$ was
arbitrary, the reduction at the start of the proof applies and $u$ is
continuously differentiable on $\Omega$.
\end{proof}
